\documentclass{article}


\usepackage{fontspec}

\setmainfont{Times New Roman}

\newfontfamily\cyrillicfonttt{Hack}

\newfontfamily\monofont{Hack}




\usepackage{needspace}



\usepackage{polyglossia}
% \setmainbreaklines=true,language{russian}


\usepackage{amsmath, amssymb, amsfonts}
\usepackage{hyperref}
\usepackage{tikz, float, pgfplots}

\usepackage{svg}


\usepackage[margin=3cm]{geometry}




\usepackage{listings}
\usepackage{color}


\usepackage{graphicx}




\usepackage{pdfpages}




\newcommand{\img}[2]{
    \begin{figure}[H]
        \centering
        \includegraphics[width=0.3\textheight]{#1}
        \caption{#2}
    \end{figure}
}



\begin{document}


%     \begin{titlepage}

%         \Large

%         \begin{center}
            
            
%             МИНИСТЕРСТВО ОБРАЗОВАНИЯ РЕСПУБЛИКИ БЕЛАРУСЬ

%             Учреждение образования «БЕЛОРУССКИЙ ГОСУДАРСТВЕННЫЙ 

%             ТЕХНОЛОГИЧЕСКИЙ УНИВЕРСИТЕТ»

%         \end{center}
        
%         \vspace{1cm}
        
%         \begin{tabular}{ll}
%             Факультет & Информационных технологий \\
%             Кафедра & Информационных систем и технологий \\
%             Специальность & 6-05-0611-01 «Информационные системы и технологии» \\
%         \end{tabular}

%         \vspace{1cm}
        
%         \begin{center}
%             {\bfseries ПОЯСНИТЕЛЬНАЯ ЗАПИСКА КУРСОВОГО ПРОЕКТА}
%         \end{center}

%         \vspace{1cm}
        
%         по дисциплине «Компьютерные языки разметки»

%         Тема: Веб-сайт «Калькулятор простых алгебраических выражений»
        
%         \vspace{2cm}
        
%         \begin{minipage}{0.5\textwidth}
%         \underline{\hspace{2cm}} \\
%         \footnotesize{Введите текст здесь}
%         \end{minipage}

% %         Исполнитель
% % студент 1 курса 1 группы		______________	Д. С. Жук
% % подпись, дата
% % Руководитель
% % 	           ассистент                           	______________		Д. В. Сазонова
% % должность, учен. степень, ученое звание	подпись, дата

%         \begin{flushright}
%             Выполнил:
            
%             Крупкевич Иван Андреевич

%             БГТУ ИСИТ 1--1
%         \end{flushright}

%         \vspace{2cm}

%         \begin{flushright}
%             Проверил:
            
%         \end{flushright}
        
        
%         \vspace{\fill}
        
%         \begin{center}
%             Минск \the\year
%         \end{center}

%     \end{titlepage}

    \includepdf[pages=1]{pz-first.pdf}


    
    \renewcommand{\contentsname}{Содержание}
    
    \tableofcontents
    \newpage
    



    \section{Постановка задачи}

    \begin{enumerate}
        \item Разработать веб-сайт с интерфейсом для ввода выражения
        и отображения результатов его вычисления.
        
        \item Написать сервер на языке программирования C, 
        который будет принимать запросы от веб-сайта, обрабатывать 
        их и возвращать результаты вычислений.
        
        \item Реализовать алгоритмы для обработки простых алгебраических выражений, 
        включая операции сложения, вычитания, умножения и деления.
    
        \item Обеспечить корректную работу веб-сайта и сервера, 
        а также обработку возможных ошибок при вводе данных пользователем.
    \end{enumerate}
    
    Сделать выводы о проделанной работе.




    \subsection{Обзор аналогичных решений}



    \subsubsection{Wolframalpha} https://wolframalpha.com

    WolframAlpha представляет собой мощный вычислительный инструмент, способный решать разнообразные математические задачи, предоставляя детальные ответы и графики.
    
    \subsubsection{Photomath} https://photomath.com/

    В отличие от WolframAlpha, приложение Photomath предлагает уникальную функцию распознавания рукописных математических выражений с помощью камеры смартфона. Этот инновационный подход позволяет пользователям с легкостью вводить математические задачи, просто сфотографировав их, и получать мгновенные ответы. Photomath фокусируется на удобстве использования и доступности, делая процесс решения математических задач более интерактивным и интуитивным для пользователей всех уровней.


    \subsection{Техническое задание}
    
    \begin{enumerate}
        \item \textbf{Информация о проекте:}
        \begin{itemize}
            \item Количество веб-страниц: 3
            \item Страницы и их содержание:
            \begin{itemize}
                \item Первая страница: содержит большую кнопку "Dive In", которая переносит пользователя в удивительный мир математики.
                \item Вторая страница: список теорем с возможностью просмотра их доказательств.
                \item Третья страница: позволяет пользователю вычислить значение математического выражения.
            \end{itemize}
        \end{itemize}
        
        \item \textbf{Формирование требований к программному продукту:}
        \begin{itemize}
            \item Разработать веб-сайт с интуитивным интерфейсом и привлекательным дизайном.
            \item Обеспечить быструю и удобную навигацию между страницами.
            \item Создать страницу с перечнем теорем и возможностью просмотра их доказательств.
            \item Разработать страницу для вычисления значений математических выражений с удобным интерфейсом для ввода данных.
        \end{itemize}
        
        \item \textbf{Формулировка задач программного продукта:}
        \begin{itemize}
            \item Создать страницу с перечнем теорем и возможностью просмотра их доказательств.
            \item Разработать функционал для вычисления значений математических выражений на отдельной странице.
            \item Обеспечить корректную работу всех элементов интерфейса и функций веб-сайта.
            \item Предоставить пользователю удобный и интерактивный опыт взаимодействия с математическими данными и вычислениями.
        \end{itemize}
    \end{enumerate}
    
    
    \subsection{Выбор средств реализации программного продукта}

    При принятии решения о выборе средств для реализации программного продукта, я решил использовать среду разработки Visual Studio Code (VS Code). VS Code предоставляет удобную и гибкую среду для написания кода, обладает широким набором расширений и инструментов, что позволяет увеличить производительность и удобство разработки. Благодаря интеграции с различными языками программирования, отладчиками и системами контроля версий, VS Code обеспечивает эффективное взаимодействие с кодом и управление проектом. В результате использования VS Code я смог повысить эффективность разработки и улучшить качество программного продукта.

    Дополнительно, для отображения математических формул на веб-сайте, я использовал библиотеку KaTeX. KaTeX обеспечивает быструю и качественную отрисовку математических выражений непосредственно в браузере, что позволяет представлять алгебраические формулы на веб-странице в удобном и профессиональном виде. Благодаря интеграции KaTeX в мой веб-сайт, пользователи могут легко вводить и просматривать математические выражения, что повышает удобство использования калькулятора и обеспечивает более качественный пользовательский опыт.

    \subsubsection{Вывод}

    Определен основной функционал программного продукта, включающий создание веб-сайта для математических вычислений с функцией перехода в интерактивный раздел, просмотром теорем и их доказательством, а также возможностью вычисления математических выражений. Для реализации проекта планируется использовать технологии KaTeX для отображения математических формул, язык программирования C для серверной части, а также HTML, CSS и JavaScript для создания пользовательского интерфейса веб-сайта. Этот стек технологий обеспечит удобство использования и функциональность программного продукта.
    




    
    \section{Проектирование страниц веб-сайта}


    \subsection{Выбор способа верстки}

    При выборе способа верстки для моего проекта было решено использовать HTML и CSS в сочетании с библиотекой KaTeX. HTML (HyperText Markup breaklines=true,language) является основным языком разметки веб-страниц, который позволяет структурировать контент и создавать элементы интерфейса. CSS (Cascading Style Sheets) используется для стилизации веб-страниц, что позволяет задавать внешний вид элементов и макет страницы.

    KaTeX представляет собой быструю и легковесную библиотеку для отрисовки математических выражений в формате TeX на веб-страницах. Благодаря своей эффективности и производительности, KaTeX обеспечивает быструю загрузку и отображение сложных математических формул без необходимости использования серверного рендеринга.

    Использование HTML, CSS и библиотеки KaTeX в проекте позволило создать удобный и функциональный веб-интерфейс для взаимодействия с математическими данными, обеспечивая пользователю понятное и привлекательное представление информации.


    \subsection{Выбор стилевого оформления}

    При выборе стилевого оформления для моего проекта было принято использовать темную тему с белым текстом и минималистичным дизайном. Темный стиль интерфейса обеспечивает комфортное восприятие контента и снижает нагрузку на глаза пользователей, особенно при длительном использовании в условиях недостаточного освещения. Белый текст на темном фоне создает контрастный эффект, делая информацию более читаемой и выделяя ключевые элементы интерфейса.

    Минималистичный дизайн был выбран для обеспечения простоты и ясности восприятия пользователем. Упрощенный стиль позволяет сосредоточить внимание на основном контенте и функционале веб-сайта, не отвлекая пользователя лишними деталями. Минимализм способствует лаконичности интерфейса, делая его более современным и эстетичным.


    \subsection{Выбор шрифтового оформления}

    Для шрифтового оформления я выбрал обычные шрифты. Основная цель была сделать текст читаемым и легко воспринимаемым. Обычные шрифты подходят для широкого круга пользователей и делают дизайн простым и доступным. Этот выбор помогает пользователям сосредоточиться на информации на сайте без лишних отвлечений.


    \subsection{Разработка логотипа}

    \begin{figure}[h]
        \centering
        \includesvg[scale=3]{logo}
        \caption{Логотип}
    \end{figure}

    При разработке логотипа была использована буква "M", стилизованная в виде числовых множеств, что символизирует связь с математикой. Название "Mather" происходит от слова "Math" - математика, что отражает основную тематику проекта.


    \subsection{Разработка пользовательских элементов}

    \subsubsection{toggle switch}
    Разработка этого элемента интерфейса включала в себя создание ползунка, который плавно перемещается между двумя состояниями, отображая выбор пользователя. Я также добавил анимацию и эффекты перехода, чтобы сделать взаимодействие с переключателем более приятным и интуитивно понятным.


    \subsection{Разработка спецэффектов}

    Для создания красивой светящейся кнопки я использовал псевдоэлемент \texttt{::before}, который добавляет дополнительный слой перед содержимым кнопки. Я задал градиентный фон с помощью \texttt{linear-gradient}, установил анимацию для изменения позиции фона, добавил размытие с помощью \texttt{filter: blur()}, и настроил переход при наведении на кнопку для изменения размытия. Это создает эффект светящейся кнопки с плавными анимациями и стилизацией.


    \subsection{Выводы}

    \begin{enumerate}
        \item \textbf{Важность пользовательского опыта}: Проектирование страниц веб-сайта играет ключевую роль в создании удобного и привлекательного пользовательского опыта. Грамотно спроектированные страницы способствуют легкому взаимодействию пользователя с сайтом и повышают его удовлетворенность.
        
        \item \textbf{Согласованность дизайна}: Важно обеспечить согласованность дизайна на всех страницах веб-сайта. Это помогает создать цельное визуальное впечатление и узнаваемый стиль, что способствует узнаваемости бренда.
        
        \item \textbf{Адаптивность и отзывчивость}: С учетом разнообразия устройств, на которых пользователи могут просматривать сайт, важно обеспечить адаптивность и отзывчивость дизайна. Это позволит сайту корректно отображаться на различных устройствах и улучшит общий пользовательский опыт.
        
        \item \textbf{Удобство навигации}: Хорошо спроектированная навигация на страницах веб-сайта помогает пользователям быстро находить нужную информацию. Четкая структура и интуитивно понятные элементы навигации улучшают пользовательскую интеракцию.
        
        \item \textbf{Тестирование и оптимизация}: Важно проводить тестирование дизайна страниц веб-сайта с целью выявления возможных проблем и улучшения пользовательского опыта. Оптимизация дизайна на основе обратной связи пользователей поможет создать более эффективный и привлекательный сайт.
    \end{enumerate}

    В целом, проектирование страниц веб-сайта играет важную роль в создании успешного онлайн присутствия. Учитывая вышеперечисленные аспекты, можно создать сайт, который не только привлечет пользователей, но и обеспечит им приятный и удобный пользовательский опыт.






    \section{Реализация структуры веб-сайта}

    Представить листинги структуры веб-сайта. Обосновать выбор элементов в коде. 
    Если это большой важный кусок кода, то покажите его в приложение в виде листинга, но укажите ссылку на приложение. Например, скрипт для создания таблиц представлен в приложении Б. В листингах интервалов после абзацев быть не должно.
    \subsection{Структура HTML-документа}
    \subsection{Добавление таблиц стилей CSS}
    \subsection{Использование стандартов XML (SVG)}
    \subsection{Управление элементами DOM}
    \subsection{Выводы}





    \section{Тестирование веб-сайта}

    \subsection{Адаптивный дизайн веб-сайта}

    \begin{itemize}
        \item \textbf{Проверка на различных устройствах:}
        Веб-сайт был протестирован на различных устройствах с разными размерами экранов, чтобы убедиться, что дизайн адаптируется корректно и отображается правильно на всех устройствах.
        
        \item \textbf{Тестирование в различных браузерах:}
        Веб-сайт был протестирован в различных популярных браузерах, таких как Google Chrome, Mozilla Firefox, Safari и Microsoft Edge, чтобы убедиться, что он отображается одинаково хорошо во всех браузерах.
        
        \item \textbf{Использование инструментов разработчика:}
        Для более детального анализа адаптивности веб-сайта были использованы инструменты разработчика браузера, позволяющие эмулировать различные устройства и размеры экранов.
        
    \end{itemize}
    
    После проведения тестирования были выявлены и исправлены возможные проблемы с адаптивностью веб-сайта, что позволило обеспечить оптимальное пользовательское взаимодействие на всех устройствах и в различных браузерах.


    \subsection{Кроссбраузерность веб-сайта}

    
    \img{index_chrome.png}{страница с калькулятором на Chrome}
    \img{getstarted_chrome.png}{начальная страница на Chrome}
    \img{theorems_chrome.png}{страница с теоремами на Chrome}
    
    \img{index_firefox.png}{страница с калькулятором на Firefox}
    \img{getstarted_firefox.png}{начальная страница на Firefox}
    \img{theorems_firefox.png}{страница с теоремами на Firefox}
    
    \img{index_mobile.png}{страница с калькулятором на Телефоне}
    \img{getstarted_mobile.png}{начальная страница на Телефоне}
    \img{theorems_mobile.png}{страница с теоремами на Телефоне}
        
    


    \subsection{Руководство пользователя}

    \img{getstarted_firefox.png}{начальная страница}
    На начальной странице представлен удобный интерфейс, позволяющий пользователям легко навигироваться по различным разделам. Нажав на кнопку "Dive In", пользователь перемещается на страницу, где представлены различные теоремы. В то же время, активация переключателя низу экрана направляет пользователя на страницу с калькулятором. Здесь пользователи могут вводить математические выражения и получать соответствующие ответы. Такой подход обеспечивает удобство использования и эффективный доступ к различным функциям веб-ресурса.


    \subsection{Выводы}

    В ходе выполнения курсового проекта был разработан веб-сайт с удобным интерфейсом, позволяющим пользователям изучать теоремы и использовать калькулятор для математических вычислений. Дополнительно к использованию стандартных технологий, был написан сервер на языке программирования C для обеспечения более эффективной обработки запросов и управления данными.


    Ссылка на GitHub: \href{https://github.com/stoopotec/Counter}{https://github.com/stoopotec/Counter}

    
    \newpage



    \section{Приложение А: листинг файла ./src/http.c}
    \lstinputlisting[frame=single,basicstyle=\ttfamily\small,breaklines=true,language=C]{../src/http.c}

    \section{Приложение Б: листинг файла ./src/http.h}
    \lstinputlisting[frame=single,basicstyle=\ttfamily\small,breaklines=true,language=C]{../src/http.h}

    \section{Приложение В: листинг файла ./src/internet.c}
    \lstinputlisting[frame=single,basicstyle=\ttfamily\small,breaklines=true,language=C]{../src/internet.c}

    \section{Приложение Г: листинг файла ./src/internet.h}
    \lstinputlisting[frame=single,basicstyle=\ttfamily\small,breaklines=true,language=C]{../src/internet.h}

    \section{Приложение Д: листинг файла ./src/list.h}
    \lstinputlisting[frame=single,basicstyle=\ttfamily\small,breaklines=true,language=C]{../src/list.h}

    \section{Приложение Е: листинг файла ./src/main.c}
    \lstinputlisting[frame=single,basicstyle=\ttfamily\small,breaklines=true,language=C]{../src/main.c}

    \section{Приложение Ж: листинг файла ./src/messages.h}
    \lstinputlisting[frame=single,basicstyle=\ttfamily\small,breaklines=true,language=C]{../src/messages.h}

    \section{Приложение И: листинг файла ./src/send\_files.c}
    \lstinputlisting[frame=single,basicstyle=\ttfamily\small,breaklines=true,language=C]{../src/send\_files.c}

    \section{Приложение К: листинг файла ./src/server.c}
    \lstinputlisting[frame=single,basicstyle=\ttfamily\small,breaklines=true,language=C]{../src/server.c}

    \section{Приложение Л: листинг файла ./src/server.h}
    \lstinputlisting[frame=single,basicstyle=\ttfamily\small,breaklines=true,language=C]{../src/server.h}

    \section{Приложение М: листинг файла ./src/transformer.c}
    \lstinputlisting[frame=single,basicstyle=\ttfamily\small,breaklines=true,language=C]{../src/transformer.c}

    \section{Приложение Н: листинг файла ./src/transformer.h}
    \lstinputlisting[frame=single,basicstyle=\ttfamily\small,breaklines=true,language=C]{../src/transformer.h}

    \section{Приложение О: листинг файла ./src/transformer\_test.c}
    \lstinputlisting[frame=single,basicstyle=\ttfamily\small,breaklines=true,language=C]{../src/transformer\_test.c}


    \section{Приложение П: листинг файла ./public/equation-list.css}
    \lstinputlisting[frame=single,basicstyle=\ttfamily\small,breaklines=true,language=C]{../public/equation-list.css}

    \section{Приложение Р: листинг файла ./public/equations.html}
    \lstinputlisting[frame=single,basicstyle=\ttfamily\small,breaklines=true,language=C]{../public/equations.html}

    \section{Приложение С: листинг файла ./public/equations.xml}
    \lstinputlisting[frame=single,basicstyle=\ttfamily\small,breaklines=true,language=C]{../public/equations.xml}

    \section{Приложение Т: листинг файла ./public/getstarted.css}
    \lstinputlisting[frame=single,basicstyle=\ttfamily\small,breaklines=true,language=C]{../public/getstarted.css}

    \section{Приложение У: листинг файла ./public/getstarted.html}
    \lstinputlisting[frame=single,basicstyle=\ttfamily\small,breaklines=true,language=C]{../public/getstarted.html}

    \section{Приложение Ф: листинг файла ./public/global.css}
    \lstinputlisting[frame=single,basicstyle=\ttfamily\small,breaklines=true,language=C]{../public/global.css}

    \section{Приложение Х: листинг файла ./public/global.js}
    \lstinputlisting[frame=single,basicstyle=\ttfamily\small,breaklines=true,language=C]{../public/global.js}

    \section{Приложение Ц: листинг файла ./public/index.css}
    \lstinputlisting[frame=single,basicstyle=\ttfamily\small,breaklines=true,language=C]{../public/index.css}

    \section{Приложение Ш: листинг файла ./public/index.html}
    \lstinputlisting[frame=single,basicstyle=\ttfamily\small,breaklines=true,language=C]{../public/index.html}

    \section{Приложение Щ: листинг файла ./public/menu.css}
    \lstinputlisting[frame=single,basicstyle=\ttfamily\small,breaklines=true,language=C]{../public/menu.css}

    \section{Приложение Э: листинг файла ./public/toggle-switch.css}
    \lstinputlisting[frame=single,basicstyle=\ttfamily\small,breaklines=true,language=C]{../public/toggle-switch.css}

    \section{Приложение Ю: листинг файла ./public/toggle-switch.js}
    \lstinputlisting[frame=single,basicstyle=\ttfamily\small,breaklines=true,language=C]{../public/toggle-switch.js}

    \section{Приложение Я: листинг файла ./public/transitions.js}
    \lstinputlisting[frame=single,basicstyle=\ttfamily\small,breaklines=true,language=C]{../public/transitions.js}

    

\end{document}